\section*{References}

Only sources actually used above, with loci sufficient to check the claims made
from them. Evidence state is given where it limits what the source can support.

\subsection*{Liturgical witnesses}

\begin{itemize}[leftmargin=1.2em, itemsep=0.15em]
\item \work{Missale Romanum}, editio typica (Vatican City: Typis Polyglottis
Vaticanis, 1962), printed pp.~397--398, marginal nos.\ 1592--1601, with p.~396
for the preceding Postcommunion and the head of p.~398 for the following
heading; and its \latin{Rubricae generales}, cited for RG~127~b, 434~b, 475~a
and 494~b. Facsimile:
\url{https://media.churchmusicassociation.org/pdf/missale62.pdf},
PDF pp.~478--479. Source-library identity
\texttt{edition.\allowbreak catholic-church.\allowbreak missale-romanum.\allowbreak vatican-typica-1962},
artifact
\texttt{artifact.\allowbreak catholic-church.\allowbreak missale-romanum.\allowbreak vatican-typica-1962.\allowbreak cmaa-facsimile-pdf}.
\textit{Controlling text for every Latin form printed above, read as page images
at 200, 400 and 500~dpi.}
\item \work{Missale Romanum}, editio iuxta typicam (New York: Benziger
Brothers, 1962), printed pp.~392--393, in the Internet Archive item
\texttt{MissaleRomanum1962RomanMissalColorLatin}. \textit{Second 1962 edition,
read on that item's page images at leaves n468 and n469 and, for locating, in
its optical text layer. Recorded as a collation of two books; nothing is
blended and no Benziger form is adopted.}
\item \work{Missale Romanum} (Ratisbon: Pustet, 1862), printed pp.~339--341, and
the page images at leaves n425 and n426 of Internet Archive item
\texttt{bub\_gb\_E7sPAAAAIAAJ}. \textit{The tracked public-domain antecedent on
which the three orations are published, and the witness that closes the
conjunction \latin{et} before \latin{rénova}. Also the witness for the pre-1955
seasonal orations at this Sunday, and an independent confirmation of
\latin{éxtrahet}.}
\ifdefined\TriptychSynthesisEdition\else
\item \work{Missale Romanum} (Vatican typical edition of 1604) and
\work{Missale Romanum} (Venice, 1570). \textit{Tracked full text. Both carry the
conjunction \latin{et} before \latin{rénova} at the Postcommunion; the 1604 is
also the witness at the Postcommunion's pointing and at the pre-1955
\latin{alia Secreta}. Optical layers only: usable for positives here and not for
negatives, the 1604 layer returning nothing for \latin{post pentecosten} while
demonstrably carrying this Sunday's whole formulary.}
\fi
\item \work{Missale Romanum}, editio typica Vaticana of 1920, and
\work{Missale Romanum}, editio IV \latin{iuxta typicam Vaticanam} (New York:
Benziger Brothers), the pre-1955 witness whose approval leaf is dated 8~December
1942. Both opened at this formulary. \textit{The two books that print
\latin{Orationes pro diversitate Temporum assignatae, ut supra} inside this
Sunday's formulary, with \latin{Aliae Secretae} and \latin{Aliae
Postcommuniones}, and that already carry the direction \latin{Præfatio de
\Sanctissima{} Trinitate.} here; they head the day \latin{semiduplex}. Both were
read in uncorrected optical layers of altar books in small type, with no page
image of either book behind them, and the layers damage numerals: the
1920 gives this Gospel's reference as \texttt{huc. 24, 2-22} and the Offertory's
as \texttt{Ps. 39, 24 et 15}, both plainly corrupt, so their agreement with the
psalm citations of the 1862 and 1962 books is an inference and not a
collation.}
\item \work{Missale Ambrosianum} (Milan: de Sirturis, 1712), the window headed
\latin{sabbato in albis missa pro baptizatis}. \textit{Eph.~3:13--21 entire with
the same liturgical incipit adaptation the Roman book makes, read as the
baptismal lesson of Easter Saturday. Optical layer only; an eighteenth-century
printing is evidence about that book and not about the medieval Ambrosian
lectionary.}
\item \texttt{src/sources/calendars/roman-1962/propers.yaml}, this repository's
tracked registry of the 1962 Mass formularies and their proper elements, parsed
as data. \textit{A repository derivative of the 1962 book's own text layer and
not a collation of its pages, and the instrument behind every statement above
that an element of this Mass is, or is not, appointed at some other formulary.
At the chants: the Gradual at the Third, Fourth, Fifth and Sixth Sundays after
the Epiphany and as the Eighteenth Sunday's Alleluia verse, the Offertory at the
Friday after Lent~II, the Communion at the Thursday after Lent~IV, and Ps.~97:1
as the Introit antiphon of the Fourth Sunday after Easter and as the psalm verse
standing at the Introits of Easter Thursday, of the Assumption, of the Maternity
of the Blessed Virgin Mary and of the Octave of the Nativity. At the Epistle: the other two
appointments of Ephesians~3, at the Sacred Heart and at St~Margaret Mary
Alacoque, and the doxology Eph.~3:20--21 standing in it at this Mass and at no
other.\ifdefined\TriptychSynthesisEdition\else{} At the Gospel: Luke~14 in three
separated blocks, with vv.~12--15 and v.~25 appointed nowhere in it.\fi{} It
holds no votive Masses at all, so the
Communion's votive appointment above rests on the 1962 book's own text layer and
not on it. No page image controls any of them.}
\ifdefined\TriptychSynthesisEdition\else
\item {\raggedright \work{The 1962 Roman Calendar: Rules, Ranks, and Reckoning}, this
repository's own reference guide, at
\texttt{src/claude/liturgy/roman-rite/1962/\allowbreak reference/\allowbreak liturgical-calendar/\allowbreak sections/\allowbreak 50-resumed-sundays.tex}.
\textit{The rule that a resumed Epiphany Sunday keeps its own orations, Epistle
and Gospel and borrows the Twenty-third Sunday's chants, printed there from
\latin{Rubricae generales} n.~18 and Missal rubric 298. Those two numbers reach
the commentary above through that guide and not through the 1962 Missal entry,
whose cited \latin{Rubricae generales} loci are the four named with it.}\par}
\fi
\item Gerbert, M., \work{Monumenta veteris liturgiae Alemannicae}, Pars~I
(St~Blasien, 1777), printed pp.~146, 178, 180 and 184, read on the page images
of two independently produced Internet Archive scans. \textit{The Frankish
Gelasian's Mass carrying all three of this Sunday's orations, with Gerbert's
per-oration sigla and his own footnotes stating the numbering arithmetic; the
September Ember Sunday at p.~178 and the \latin{Dominica vacat} printed as a Mass
at p.~180. The first scan clips the outer margin of p.~184, so quoted line-ends
come from the second. The page images establish what Gerbert printed, not what
any manuscript reads.}
\item Sacred Congregation of Rites, general decree \work{De rubricis ad
simpliciorem formam redigendis}, 23~March 1955, in \work{Acta Apostolicae Sedis}
47 (1955) 218--224, Tit.~II n.~1, Tit.~V~a) n.~1 and Tit.~V~b) n.~7.
\textit{Read in the Holy See's own optical transcription of its own \work{Acta}
at vatican.va, and at 200~dpi of that transcription rather than of a facsimile.
One witness only: no printing of 1955 was collated, and neither the decree nor
this volume of the \work{Acta} is registered in this repository's sources.}
\item Brightman, F.~E., \work{The English Rite}, vol.~II (London: Rivingtons,
1915), printed pp.~520--521; and Blunt, J.~H., \work{The Annotated Book of
Common Prayer} (London, 1866), at the Sixteenth and Seventeenth Sundays after
Trinity. \textit{The two synopses that carry the Sarum arrangement of this
formulary into the English books of 1549, 1552 and 1661. Optical layers, no page
image, and no printing of the Book of Common Prayer itself was collated; they
tag the Collect \latin{(Greg. 172)} and \latin{Greg. Orationes Quotidianae}
respectively, neither expanded.}
\item Wilson, H.~A., ed., \work{The Gelasian Sacramentary} (Oxford: Clarendon,
1894). Book~III sect.~XII no.~693, printed p.~231 (the Secret and the
Postcommunion, under \latin{ITEM ALIA MISSA}); no.~692, p.~230 (the Fifteenth
Sunday's oration set); the Appendix at p.~358 for \latin{Gerb. Hebd. xx post
Pentecosten} and for \latin{Gerb. Dominica Vacat}; and Wilson's footnote to
no.~693 for Pamelius. \textit{Read in an uncorrected optical layer; no page image
opened; the numbering is Wilson's own.}
\item Wilson, H.~A., ed., \work{The Gregorian Sacramentary under Charles the
Great}, Henry Bradshaw Society XLIX (1915), sect.~XXXV, printed p.~174, with
sects.~XXXIII and XXXIIII for the two preceding Sundays and the marginal
antiphonary cues of Cod.~Ottobonianus lat.~313 at sect.~XXXIII, p.~173.
\textit{Optical layer only. Where the cue block attaches is disputed and no page
image was opened.}
\item \work{Sacramentarium Leonianum}, ed. C.~L. Feltoe (Cambridge, 1896),
searched whole. \textit{Feltoe's own title for the sacramentary the prose above
calls the Veronense. None of this Sunday's three orations is in it; the one
near miss, \latin{tua nos semper gratia praeveniens laetificatur}, is a
different prayer.}
\item Dickinson, F.~H., ed., \work{Missale ad usum insignis et præclaræ
ecclesiæ Sarum} (Burntisland, 1861--1883), printed cols.\ 506, 509--510, 512.
\textit{The split of this formulary across Trinity XVI and XVII, and the two
Sarum alleluias that differ from the Roman. Read in an uncorrected optical layer
of that scan; no page image opened.}
\item \work{Liber Comitis}, in J.-P. Migne, \work{Patrologia Latina} 30, col.\
522, read on the page image of leaf n264, one of seven page images of leaves
n261--n267 of the Internet Archive digitisation. \textit{The Gospel under \latin{Dominica mensis viii} paired with
Ephesians 4:1--6; the headings read at high zoom because the optical text
mangles the numerals.}
\item Schuster, I., \work{The Sacramentary (Liber Sacramentorum)}, vol.~III,
English of Arthur Levelis-Marke (London: Burns Oates \& Washbourne, 1927),
printed pp.~142--144. \textit{The older title \latin{Prima post natale Sancti
Cypriani}, the stations of the surrounding days, the course of Ephesians, and
three of the four second appointments of the chants. Read in an optical layer of
the registered item; no page image opened, and the title is a reported title and
not a collated manuscript reading.}
\item Guéranger, P., \work{The Liturgical Year}, vol.~XI, English of Dom
Laurence Shepherd (Dublin: Duffy, 1900), the chapter for this Sunday at printed
pp.~356--371,%
\ifdefined\TriptychSynthesisEdition{} quoted above at p.~357 (the Collect), at
one place within pp.~358--360 (the four measures), at p.~368 (Ecclus.~3:20--21
joined to the Gospel's closing axiom), at p.~370 (the Secret) and at p.~371 (the
Postcommunion), with his readings of the Gradual at p.~363, of the parable's
wedding at p.~365, of Ambrose as the homilist the Church has chosen for this
Sunday at p.~367, of the Offertory at p.~370 and of the Communion at
pp.~370--371 reported and not quoted. \textit{The chapter heading, the printed folio
and the opening sentence were read on a rendered page image at printed p.~356,
and pp.~357, 359, 363, 365, 367--368 and 370--371 were rendered likewise; two
passages of the chapter, at printed pp.~358 and 360, rest on the volume's own
extracted text layer, with no page image behind them. A
nineteenth-century liturgist writing devotionally, cited as that; a search of the
whole chapter for the vocabulary of liturgical history returns nothing.}%
\else{} quoted above at pp.~357 (the Collect), 358--360 (the Epistle and
the four measures), 363 (the Gradual), 365 and 367--368 (the Gospel, the wedding
and Ecclus.~3:20--21), 370 (the Offertory and the Secret) and 370--371 (the
Communion and the Postcommunion). \textit{The chapter heading, the printed folio
and the opening sentence were read on a rendered page image at printed p.~356,
and pp.~357, 359, 363, 365, 367--368 and 370--371 were rendered likewise; the
sentence quoted from p.~360 and the observation reported from p.~358 rest on the
volume's own extracted text layer, with no page image behind them. A
nineteenth-century liturgist writing devotionally, cited as that; a search of the
whole chapter for the vocabulary of liturgical history returns nothing.}%
\fi
\item Fortescue, A., ``Preface,'' in \work{The Catholic Encyclopedia}, vol.~XII
(New York: Robert Appleton, 1911); and Fortescue, A., \work{The Mass: A Study of
the Roman Liturgy} (London: Longmans, 1922), searched whole. \textit{The
encyclopedia article carries the 1759 attribution of the Sunday Trinity preface
to Clement~XIII and names no \work{Acta} locus; the book by the same author
carries neither the rule nor the date, so the two do not corroborate each other.}
\item gregorien.info, the chant database's index of R.-J. Hesbert,
\work{Antiphonale Missarum Sextuplex}, cited by that database in its 1985
reprint: the day pages for AMS 188 and its neighbours, five chant records, six
per-manuscript indexes and the manuscript index. \textit{The 1985 reprint is
the printing behind every AMS number reported above, and it is the book a
reader checks them in. Hesbert's own 1935 Vromant printing is a different
object: the source records of this repository carry a catalog entry for it and
no artifact, and no page of it stands behind anything here. Every AMS result
above is this database's report of Hesbert's numbering and not his page. The
siglum key was reconstructed by cross-tabulation and is not printed by the
site; the whole-calendar view is capped by the host at 500 rows.}
\ifdefined\TriptychSynthesisEdition\else
\item Lasance, F.~X., \work{The New Roman Missal} (New York: Benziger, 1945),
optical layer. \textit{Corroborates at optical state that the Thursday after
Lent~IV carries this Sunday's Communion and the Fifteenth Sunday's Gospel.}
\fi
\end{itemize}

\subsection*{Scripture and its English}

\begin{itemize}[leftmargin=1.2em, itemsep=0.15em]
\item \work{Vulgata Clementina}, the Clementine Vulgate, tracked at
\texttt{src/sources/bibles/clementine-vulgate}. \textit{The comparison text for
every divergence
\ifdefined\TriptychSynthesisEdition recorded\else tabulated\fi{} above, and the
corpus behind the exhaustive lists of
\latin{mitis}, \latin{canticum novum} and \latin{Amice}.}
\item \work{The Holy Bible, Douay--Rheims}, revised by Bishop Richard Challoner
(1749--1752), Project Gutenberg e-text, tracked at
\texttt{src/sources/bibles/douay-rheims}.%
\ifdefined\TriptychSynthesisEdition{} \textit{The registered public-domain
English of the seven scriptural elements, and the English this edition quotes
its scriptural wording from. The three per-book artifacts it is quoted from, at
Psalms, Ephesians and Luke, are named in
\texttt{research/source-bindings.toml}.}%
\else{} \textit{The registered public-domain
English of all seven scriptural elements, printed whole where a chant sings part
of a verse. The three per-book artifacts it is quoted from, at Psalms,
Ephesians and Luke, are named in \texttt{research/source-bindings.toml}.}%
\fi
\item \work{The Roman Missal translated into the English language for the use of
the laity}, first revised edition (Philadelphia: Eugene Cummiskey, 1861),
formulary \latin{XVI. SUNDAY after PENTECOST}, printed pp.~429 and 431, read on
the page images at leaves n437 and n439 of Internet Archive item
\texttt{romanmissaltran00churgoog}; passage
\texttt{passage.\allowbreak eugene-cummiskey.\allowbreak roman-missal-english-laity.\allowbreak philadelphia-1861.\allowbreak post-pentecosten-16}.
\ifdefined\TriptychSynthesisEdition
\textit{The registered public-domain English of the three orations, a free
devotional rendering this edition reports on and does not reproduce. Its English
of the scriptural elements is not used.}%
\else
\textit{The registered public-domain English of the three orations, a free
devotional rendering reproduced with the book's own spelling. Its English of the
scriptural elements is not used.}%
\fi
\item \work{The Authorized Version of the Bible}, Project Gutenberg eBook 10.
\textit{The English the gallery's later users quote, and the corpus in which
``inner man'' and ``lowest room'' were counted.}
\ifdefined\TriptychSynthesisEdition\else
\item \work{The Revised Version} (1895), tracked here. \textit{The English
against which the shorter Latin clause at Luke~14:3 was noticed; no Greek New
Testament is tracked, so the manuscript question is not settled.}
\fi
\end{itemize}

\subsection*{Patristic, medieval and conciliar witnesses}

\begin{itemize}[leftmargin=1.2em, itemsep=0.15em]
\item Augustine, \work{Enarrationes in Psalmos}: on Ps.~85 \S\S1--3, 5, 7; on
Ps.~101 sermo~I \S\S16--19; on Ps.~97 \S1; on Ps.~39 \S\S23--25; on Ps.~70
sermo~II \S\S1--5. Latin from the Migne transcriptions served by la.wikisource
and from the augustinus.it texts; English from the \work{Nicene and Post-Nicene
Fathers}, first series, vol.~8, as New Advent delivers it. The fifteen exact
artifacts, three per psalm, are named in
\texttt{research/source-bindings.toml}. \textit{Web transcriptions of Migne, not critical editions;
cited by enarratio, sermo and section because no Migne column was verified
against a facsimile.}
\item Augustine, \work{Enchiridion ad Laurentium} 32; \work{De gratia et libero
arbitrio} 17.33. \textit{The Collect's doctrine stated outside any commentary on
the Collect.}
\item Augustine, \work{Sermo} 165, \S\S1.1, 2.2, 3.3--5.5, 7.9; \work{Epistula}
55.14.25; \work{In Iohannis euangelium tractatus} 118.5. \textit{The three
cross-readings of Eph.~3:18, with three different sets of virtues; and the
defence of \latin{altitudo} against the missal's \latin{sublímitas}.}
\item Augustine, \work{Quaestiones euangeliorum} II q.~29. \textit{The dropsical
man compared to the animal in the well and to the covetous rich man. The
retrieved transcription carries optical damage in that paragraph, so the sense
is reported above and the wording is not quoted.}
\item Ambrose, \work{Expositio evangelii secundum Lucam} VII.195--196, with
books VIII and IX read whole for the negative, in the same Migne transcription
served by la.wikisource.
\item Bede, \work{In Lucae euangelium expositio}, PL~92 cols.\ 510D--513A, with
the prefatory letter to Acca. \textit{The fullest continuous Latin exposition of
this pericope. Bede states there that he marked his borrowings with their
authors' initials; the transcription read here carries no such marks.}
\item John Chrysostom, \work{Homilies on the Epistle to the Ephesians}, Homily~7
(with 6 and 8 as boundary witnesses), in the \work{Nicene and Post-Nicene
Fathers}, first series, vol.~13, as New Advent delivers it.
\textit{Read in English; the Greek of PG~62 was not opened, so no argument above
rests on his wording.}
\item Cyril of Alexandria, \work{A Commentary on the Gospel according to
S.~Luke}, Sermons 101 and 102, English of R.~Payne Smith (Oxford, 1859),
pp.~471--479. \textit{The commentary survives entire only in Syriac and no Greek
or Syriac witness was opened, so Cyril's argument is used above and his wording
is not quoted. His lemma at Luke~14:5 reads ``son \dots\ or ox'' and does not
answer the missal's \latin{ásinus}.}
\item Benedict of Nursia, \work{Regula} 7, \latin{De humilitate}, opening
sentences. \textit{Quotes the Gospel's last verse as the voice of Scripture and
names no chapter.}
\ifdefined\TriptychSynthesisEdition\else
\item Gregory the Great, \work{Regula pastoralis} III.18 and III.21, in the
\work{Nicene and Post-Nicene Fathers}, second series, vol.~12. \textit{Its
editors refer the sentence at III.18 to Luke~18:14, and Gregory's nearest
citation of Luke~14 is at III.21 and is of vv.~12--14, which the Missal does not
appoint.}
\fi
\item Thomas Aquinas, \work{Super Epistolam ad Ephesios lectura} cap.~3
\ifdefined\TriptychSynthesisEdition lect.~5\else lect.~4--5\fi{}, Corpus
Thomisticum, Marietti text; \work{Catena aurea in Lucam}
cap.~14 lect.~1--2, Turin 1953; \work{Postilla super Psalmos} in ps.~39
nn.~6--7. \textit{The Catena names eight authorities at this pericope, of which
four --- Theophylact, Gregory's \work{Moralia}, Basil and a Chrysostom excerpt
--- were verified at no locus and are not quoted above.}
\item Cassiodorus, \work{Expositio psalmorum}, in Ps.~XXXIX \S\S21--24; in
Ps.~LXX \S\S4, 21--23, 32--33; in Ps.~LXXXV \S\S2, 4, 6, 8, 10, 18, 23--24; in
Ps.~XCVII \S\S2, 4, 6--7, 14; in Ps.~CI \S\S2--3, 22--23. Read in the Corpus
Corporum text as monumenta.ch delivers it, per psalm. \textit{A Migne-derived web
transcription, not Adriaen's CCSL~98 and not PL~70 at a column; the section
numbers are the host's, and the loci above are cited by lemma because the page's
verse numerals run three ahead of the missal at Ps.~39 and one ahead at Ps.~97.
His lemma stands with the missal against the Clementine at four of the five
chants, and no critical text was collated at those four lemmata.}
\item Theodoret of Cyrus, \work{Interpretatio in Psalmos}, in J.-P. Migne,
\work{Patrologia Graeca} 80: cols.\ 1159--1160 (Ps.~39), 1423--1426 (Ps.~70),
1553--1556 (Ps.~85), 1657--1658 (Ps.~97) and 1679--1682 (Ps.~101).
\textit{Migne's reprint of Schulze's Halle recension with a facing Latin version,
not a critical Greek text; every quotation above is Migne's Latin, his Greek
named where the two part and quoted nowhere, and the two differ in tense at
Ps.~101:17. Cols.\ 1423--1424 and 1553--1554 --- the openings of his Ps.~70 and
his Ps.~85 --- were read on the facsimile's uncorrected optical layer; the other
six ranges, cols.\ 1159--1160, 1425--1426, 1555--1556, 1657--1658, 1679--1680
and 1681--1682, on rendered page images. The stable locus in this facsimile is
the column and not the artifact's page, pp.~878--879 reprinting the opening of
pp.~876--877. R.~C. Hill's English is in copyright and was not used.}
\item Bellarmine, R., \work{Explanatio in Psalmos}, in the English of John
O'Sullivan as \work{A Commentary on the Book of Psalms} (Dublin: Duffy, 1866),
printed pp.~122 (Ps.~39), 212--213 (Ps.~70), 259--261 (Ps.~85), 306 (Ps.~97) and
317 (Ps.~101). \textit{Reception in English only: the translator declares the
abridgement, so nothing above is Bellarmine's Latin and no argument rests on an
absence from these pages. Printed pp.~212--213, 259 and 306 (artifact PDF
pp.~222--223, 269 and 316) were read on the facsimile's uncorrected optical
layer; printed pp.~122, 260 and 317 (PDF pp.~132, 270 and 327) on rendered page
images; for printed p.~261, the tail of the Ps.~85 span, no state is
recorded. His printed verse numbers run one below the Vulgate at
Ps.~39 and Ps.~101 and level at the other three, and the argument line and psalm
text printed above each exposition are the Douay--Rheims and not his.}
\ifdefined\TriptychSynthesisEdition
\item Jerome, \work{Epistula} 106 (\latin{Ad Sunniam et Fretelam}) \S23.
\textit{His ruling for \latin{respice} against a reported \latin{festina} at
the Offertory's clause. The
registered witness is a 1937 English translation whose copyright the source
record marks unestablished: the Latin lemmata and the words at issue are quoted
above and the translator's English is not, and the Latin original of the letter
was not opened.}
\else
\item Jerome, \work{Epistula} 106 (\latin{Ad Sunniam et Fretelam}) \S\S23 and
43. \textit{His ruling for \latin{respice} against a reported \latin{festina} at
the Offertory's clause, and his two rulings on the Communion's wording. The
registered witness is a 1937 English translation whose copyright the source
record marks unestablished: the Latin lemmata and the words at issue are quoted
above and the translator's English is not, and the Latin original of the letter
was not opened.}
\fi
\ifdefined\TriptychSynthesisEdition\else
\item Cornelius a Lapide, \work{Commentaria in omnes D. Pauli Epistolas}
(Antwerp, 1614), in the Internet Archive facsimile of that printing, whose
optical layer of 97,832 lines was searched whole. \textit{Used above for the
printing's extent and for nothing else. Its last running heads are
\latin{COMMENT. IN EPIST. AD GALAT.}, no \latin{ad Ephesios} running head stands
anywhere in it, and the 54 hits on \latin{Ephes} variants are cross-references
inside commentary on other epistles. An uncorrected optical layer; no page image
was opened.}
\fi
\item Council of Trent,
\ifdefined\TriptychSynthesisEdition
Session~VI, cap.~X and cap.~XVI, from the
\else
Session~VI, cap.~X, cap.~XVI and can.~32, from the
\fi
Tauchnitz 1887 printing of the \work{Canones et decreta}, page-verified in this
repository. \textit{Cap.~X quotes a Sunday collect of this
series, which the printing's own footnote identifies as the Thirteenth Sunday's,
and quotes it altered.}
\end{itemize}

\subsection*{Sources of the page-2 chronology}

\begin{itemize}[leftmargin=1.2em, itemsep=0.15em]
\item \work{The Catholic Encyclopedia} (New York: Robert Appleton, 1908--1912),
identified by volume and by New Advent article number, which is what the
Scripture chronology corpus records: vol.~V, article 05485a and vol.~XI, article
11567b for the Epistle's two labels; vol.~XIV, article 14530a for the Gospel's;
vol.~III, articles 03738a and 03731a for the two Nativity labels; and vol.~IV,
article 04642b for the reign four of the psalm dossiers attribute to David.
\textit{Cited only for the labels the corpus takes from them, and no article
title is assigned to any of them here because the corpus records none; none was
opened for a place of composition, a first audience or a life stage.}
\item Pontifical Biblical Commission, \work{De auctore, tempore et veritate
evangeliorum Marci et Lucae}, responsa, \work{Acta Apostolicae Sedis} 4 (1912).
\textit{Fixes the writing of Luke relative to Acts and to the Roman imprisonment
rather than by a year.}
\item \work{New American Bible Revised Edition}, the introduction to the Psalms,
United States Conference of Catholic Bishops. \textit{The source of the
whole-Psalter composition boundary the corpus returns for all five appointed
psalms. Its bytes are registered restricted and nothing of it is reproduced
here: what the Date column prints for that boundary is the generated projection
of the corpus's structured date, not the introduction's own wording.}
\item Haydock, G.~L., \work{Douay Rheims Bible with Haydock Commentary}, at
Ps.~70. \textit{The Ussher-attributed marginal year of the world the corpus
carries as a reported figure, excluded from its own basis.}
\end{itemize}

\subsection*{Later users, in the gallery}

\begin{itemize}[leftmargin=1.2em, itemsep=0.15em]
\item Jerome, Jerome K., \work{Three Men in a Boat (To Say Nothing of the Dog)}
(London: J.~W. Arrowsmith, 1889), chapter XIX. Project Gutenberg eBook 308.
\item Trollope, Anthony, \work{Barchester Towers} (London: Longman, Brown,
Green, Longmans, \& Roberts, 1857), chapter VII. Project Gutenberg eBook 3409.
\item \work{San Antonio Daily Light} (San Antonio, Tex.), 7 October 1897, p.~8,
``THE EMERGENCY MET''; and \work{Evening Star} (Washington, D.C.), 26 September
1942, image 21, set aside as a sermon title. Library of Congress,
\work{Chronicling America}; the phrase search for ``ox is in the ditch'' returns
324 pages, of which the first page of ten records was retrieved.
\item Johnson, Clifton, \work{Highways and Byways of the South} (New York: The
Macmillan Company, 1904), p.~169.
\item Rossetti, Christina G., ``Sit Down in the Lowest Room.'',
\work{Macmillan's Magazine}, no.~53, vol.~IX (London: Macmillan, March 1864),
pp.~436--439; and \work{The Poetical Works of Christina Georgina Rossetti}, ed.
William Michael Rossetti (London: Macmillan, 1904), pp.~16--19 with the editor's
note at pp.~460--461.
\item Barrett Browning, Elizabeth, \work{Sonnets from the Portuguese}, Sonnet
XLIII, in \work{Poems} (London: Chapman and Hall, 1850). Project Gutenberg
eBook 2002.
\item Phillipson, John S., ``\,`How Do I Love Thee?' --- an Echo of St.~Paul,''
\work{The Victorian Newsletter}, no.~22 (Fall 1962), p.~22, note~7. \textit{In
copyright; served openly by the Ohio State University Libraries and quoted short
and attributed.}
\item Nietzsche, Friedrich, \work{Menschliches, Allzumenschliches: Ein Buch für
freie Geister}, Erster Band, Zweites Hauptstück, aphorism 87, Project Gutenberg
eBook 7207; and \work{Human, All-Too-Human: A Book for Free Spirits}, Part~I,
translated by Helen Zimmern (Edinburgh and London: T.~N. Foulis, 1910), printed
p.~88, the page fixed by the running heads that bracket it. \textit{The aphorism
is not in \work{Jenseits von Gut und Böse}, \work{Also sprach Zarathustra},
\work{Götzen-Dämmerung}, \work{Ecce homo} or \work{Der Wille zur Macht}, all
searched.}
\item Herbert, Henry William, ``The Yorkshire Moors,'' in \work{Life and
Writings of Frank Forester}, ed.\ David W. Judd, vol.~2 (New York: Orange Judd
Company, 1882), p.~35, the page fixed by its running head. \textit{The gastric
sense of ``inner man'' in a checked witness of 1882, seven years before Jerome
K.~Jerome.}
\item Cervantes Saavedra, Miguel de, \work{Don Quijote}, Primera parte,
cap.~XI, Project Gutenberg eBook 2000; and John Ormsby's English, Project
Gutenberg eBook 996. \textit{The Spanish proverbial form of Luke~14:11 set aside
in the appendix above, read in those two e-texts and in no printed edition.}
% The last entry of the References is guarded, because what follows this file
% is the terminal apparatus: main.tex inputs generation-metadata.tex, which
% prints the revision-timestamp line, and then issues the reuse-and-rights
% colophon.  guidance/editorial.md and guidance/repository.md both require that
% block to be read under the terminal content on one physical final page and
% forbid it to stand on a page of its own, so where it cannot follow this entry
% the break has to fall BEFORE the entry and carry the entry down with it.
% Measured on this run's 200 dpi review raster: this entry sets to four lines
% (132 px), the list end and the timestamp line take 43 px and 33 px after it,
% and main.tex's own \Needspace{3\baselineskip} is then tested against the page
% goal before \TriptychRightsNotice recovers the footer, so it asks a further
% 99 px.  That is 307 px, or 9.3 baselines; ten are demanded.  The colophon's
% own block is paid for by the 116 px (0.58 in) of suppressed footer the notice
% recovers and needs nothing here.  The two editions break this shared text in
% different places and the guard answers each on its own terms: the canonical
% References reach to within about five baselines of the foot of their last
% page, so it fires and this entry goes down to the final page ahead of the
% colophon; the companion's end about twenty-nine baselines clear of the foot,
% so it passes and nothing in that edition moves.
\Needspace{10\baselineskip}%
\item Negative corpora for the gallery: E.~C. Brewer, \work{Dictionary of Phrase
and Fable}; John Camden Hotten, \work{The Slang Dictionary}; Francis Grose,
\work{A Classical Dictionary of the Vulgar Tongue}; Webster's Unabridged of
1913; and a ninety-volume Project Gutenberg sample of English and American
literature, including seven novels of Dickens and Scott's \work{Waverley} and
\work{Rob Roy}. \textit{Searched without result for ``inner man''.}
\end{itemize}
